% !TEX root = main.tex
\section{Distribution of Collatz Primitive Source}~\label{sec:distribution-of-collatz-primitive-source}

Now let's consider the Collatz primitive source mentioned in \cref{thm:layers-of-reduce-space-and-source-value-set}. They are those number that can't trace back again.
And compaired with the total primitive source, we are more concerned about its distribution in source value set. 

\begin{definition}[Collatz primitive source]~\label{collatz-primitive-source}
    % For $b\in \widehat{\mathbb{N}}_{\beta}$, if there is a sequence $\bm{v}\in \mathbb{N}^{n}$ such that $f_{n}(\bm{v})$ is not reversible modulo $\alpha$, or
    % \[\Log_{\beta, \alpha}(f_{n}(\bm{v})^{-1} \cdot \gamma) = \varnothing, \]
    % then we say $\bm{v}$ is a $n$-step Collatz primitive source vector of $b$ and $f_{n}(\bm{v})$ is a $n$-step Collatz primitive source of $b$.
    For $b\in \widehat{\mathbb{N}}_{\beta}$, if not exists $m \ge 1$ and $a \in \widehat{\mathbb{N}}_{\beta}$ such that
    \[ b \cdot \beta^{m} = \alpha \cdot a + \gamma, \]
    then we say $b$ is a Collatz primitive source.
    All primitive source are denoted as $\mathbb{S}$.
    If $a \not\in\mathbb{S}$, we say $a$ is Collatz reversible.
\end{definition}
\begin{proposition}
    If $\alpha \mid \gamma$, then $\mathbb{S} = \{b \in \widehat{\mathbb{N}}_\beta: a\nmid b\}$.
\end{proposition}
\begin{proof}
    If $\alpha\mid \gamma$, $\forall a\in\widehat{\mathbb{N}}_\beta$, we have
    \[ a_1 \cdot \beta^{m} = \alpha \cdot a + \gamma = \alpha \cdot \left(a + \frac{\gamma}{\alpha}\right). \]
    That means $\alpha\mid a_1$. Thus if $\alpha \nmid b$, then $b$ is a primitive source.
\end{proof}
From \cref{cor:modulo-alpha-of-layers}, we have
\begin{proposition}
    If $\alpha \mid \gamma$, $\mathbb{F}_{1}(b) \not\subset \mathbb{S}$.
\end{proposition}
Suppose $\alpha \nmid \gamma$.
Since if $b \cdot \beta^m = \alpha\cdot a + \gamma$, then $m \in \Log_{\beta,\alpha}(b^{-1}\cdot\gamma)$.
But not all $b$ is reversible modulo $\alpha$, and not all $b^{-1} \cdot \gamma$ is a remainder of a power of $\beta$.
\begin{definition}
    Suppose $\alpha \nmid \gamma$. Given $b \in \mathbb{S}$. Or it is not reversible modulo $\alpha$, or
    \[\Log_{\beta, \alpha}(b^{-1} \cdot \gamma) = \varnothing. \]
    Notice that $b$ is not reversible modulo $\alpha$ equivalents $\gcd(b, \alpha) > 1$. For these, we say they are trivial primitive source.
    Others we call them non-trivial primitive source.
    And all trivial primitive source are denoted as 
    \[\hat{\mathbb{S}} = \{b \in \widehat{\mathbb{N}}_{\beta} \mid \gcd(b, \alpha) > 1\} = \bigcup_{p\in\Pi(\mathbb{P}\mid\alpha)} p\widehat{\mathbb{N}}_{{\beta}}, \]
    all non-trivial primitive source are denoted as
    \[ \tilde{\mathbb{S}} = \{b \in \widehat{\mathbb{N}}_{\beta} \mid \gcd(b, \alpha) = 1 \text{ and } \Log(b^{-1}\cdot\gamma) = \varnothing\}. \]
\end{definition}
Immidiately, we have
\begin{proposition}
    $\mathbb{S} = \tilde{\mathbb{S}} \cup \hat{\mathbb{S}}$.
\end{proposition}
\begin{proposition}~\label{prop:primitive-source-of-prime-alpha}
    If $\alpha \in \mathbb{P}$ and $\ord_{\alpha}\beta = \alpha - 1$, then $\tilde{\mathbb{S}} = \varnothing$.
\end{proposition}
Suppose $\alpha \nmid \gamma$. From \cref{cor:modulo-alpha-of-layers}, we have
\begin{proposition}
    If $\mathbb{F}_{1}(b) \cap \mathbb{S} \ne \varnothing$ and $\alpha^{2}\mid (\beta^{\Phi_1} - 1)$, then $\mathbb{F}_{1}(b) \subset \mathbb{S}$.
\end{proposition}
\begin{proposition}
    If $\mathbb{F}_{1}(b) \subset \mathbb{S}$, then $\exists d \in \Pi(\mathbb{N}\mid\alpha)$ that $d > 1$, such that $d\alpha\mid \gamma (\beta^{\Phi_1} - 1)$.
\end{proposition}
\begin{proof}
    Let $d = \gcd(P\cdot\gamma, \alpha)$.
    At first we have $A_k \equiv kP\cdot\gamma + A_0 \pmod{\alpha} \in \mathbb{S}$. If $d = 1$ then $A_k$ runs all remainder modulo $\alpha$. From \cref{thm:primitive-source-are-periodic}, we will derive that $\mathbb{N}\subset\mathbb{S}$. That's impossiable. Thus $d > 1$. Let
    \[ d = \gcd(P\cdot\gamma, \alpha) \in \Pi(\mathbb{N}\mid\alpha), \]
    we have $d \mid P\cdot\gamma = \frac{\beta^{\Phi_1} - 1}{\alpha}\cdot \gamma$, i.e. $d\alpha\mid \gamma (\beta^{\Phi_1} - 1)$.
\end{proof}
And from \cref{cor:modulo-alpha-of-layers}, we also have
\begin{theorem}~\label{thm:source-set-dijoint-primitive-source}
    Suppose $\alpha\in\mathbb{P}$, $\alpha^{2}\nmid (\beta^{\Phi_1} - 1)$ and $\alpha\nmid\gamma$.
    For $b\not\in\mathbb{S}$, we have $\forall n\ge 1$, $\mathbb{F}_{n}(b) \not\subset \mathbb{S}$ and $\mathbb{F}_{n}(b) \cap \mathbb{S} \ne \varnothing$ and $\forall a \rightsquigarrow b$, $\mathbb{F}_{1}(a)/\alpha\mathbb{Z} = \mathbb{Z}/\alpha\mathbb{Z}$.
\end{theorem}
\begin{lemma}~\label{lem:all-source-value-set-has-primitive-source}
    For $b \not\in\mathbb{S}$, if there is $\mathbb{F}_{1}(b) \not\subset \mathbb{S}$, then we have $\forall n\ge 1$, $\mathbb{F}_{n}(b) \not\subset \mathbb{S}$ and $\mathbb{F}_{n}(b) \cap \mathbb{S} \ne \varnothing$.
\end{lemma}
We can  verify that arguments $(\alpha, \beta, \gamma) = (3,2,1)$ follow the requirements of \cref{thm:source-set-dijoint-primitive-source}. This is precisely the type of Collatz system we will be discussing.

Finally, at end of this ection, we provide a theorem to describe the distribution of primitive sources in $\widehat{\mathbb{N}}_{\beta}$:
\begin{theorem}~\label{thm:primitive-source-are-periodic}
    Let $b \in \mathbb{S}$. $\forall m \in \mathbb{Z}$, if $b + m\alpha > 0$, then $\ram_{\beta}(b + m\alpha) \in \mathbb{S}$.
\end{theorem}
\begin{proof}
    Trivial for $\alpha \mid \gamma$. Suppose $\alpha \nmid \gamma$.

    For those $b \in \hat{\mathbb{S}}$, we have $\gcd(b, \alpha) > 1$. Let $d = \gcd(b, \alpha)$, then we have $\gcd(d, \beta) = 1$. Thus $d\mid \ram_{\beta}(b + m\alpha)$ means $\ram_{\beta}(b + m\alpha) \in \hat{\mathbb{S}}$.

    For those $b \in \tilde{\mathbb{S}}$, we have $\gcd(b + m\alpha,\ \alpha) = 1$. If $\ram_{\beta}(b + m\alpha) \not\in \mathbb{S}$, then $\exists k \ge 1$, such that $\beta^{k}\equiv \ram_{\beta}(b + m\alpha)^{-1} \cdot \gamma \pmod{a}$. Thus we have
    \vspace{-0.5em}
    \begin{align*}
        \beta^{k - \fac_{\beta}(b + m\alpha)} &= \beta^{k}\cdot \raf_{\beta}(b + m\alpha)^{-1} \\
        &\equiv \raf_{\beta}(b + m\alpha) \cdot \ram_{\beta}(b + m\alpha)^{-1} \cdot \gamma &&\pmod{a} \\
        &\equiv (b + m\alpha)^{-1} \cdot \gamma &&\pmod{\alpha} \\
        &\equiv b^{-1} \cdot \gamma &&\pmod{\alpha},
    \end{align*}
    which implies that $b \not\in \mathbb{S}$, conflicts with the hypothesis. Therefore, it indicates that there must be $\ram_{\beta}(b + m\alpha)\in \tilde{\mathbb{S}}$.
\end{proof}
\begin{remark}
    Notice that $\beta^{k}$ are periodic modulo $\alpha$, here we should pick a large enough $k$ to make sure $k - \fac_{\beta}(b + m\alpha) \ge 0$.
\end{remark}
In fact, not only the original source, but all source values that preserve the reversible structure (reduction sequence) follow a rule similar to \cref{thm:primitive-source-are-periodic}. This is why we discuss $\Phi_{k+1} \mid \alpha\Phi_{k}$, where these source values that preserve the reversible structure are distributed at intervals of some factor of $\alpha$. If we expand our perspective from $\mathbb{N}$ to $\mathbb{Q}$ but those denominator is a power of $\alpha$, then we can find that $b + m\alpha$ also preserves a reduction sequence centered at $b$.